\begin{figure}[htbp]
\centering
\begingroup
\setlength{\tabcolsep}{0pt}
\renewcommand{\arraystretch}{1.08}
\begin{tabular}{c}
\fbox{\parbox{.88\linewidth}{\centering
\textbf{Arithmetic relations and chain polynomials}\\
Equal spectra modulo $2$; ordinary orders $2$ and $3$.}}\\[3pt]
$\Downarrow$\\[3pt]
\fbox{\parbox{.88\linewidth}{\centering
\textbf{Two separate seed identifications}\\
Real Seifert forms \quad and \quad integral monodromy modules.}}\\[3pt]
$\Downarrow\quad$ common fifth and seventh powers;\\
\textit{symmetrize and take one tensor square}\\[3pt]
\fbox{\parbox{.88\linewidth}{\centering
\textbf{Equivariant isometries over $\Q$ and every $\Z_\ell$}\\
A second tensor square cancels their eigenspace determinant defects.}}\\[3pt]
$\Downarrow\quad$ common cubic, eleventh, and thirteenth powers\\[3pt]
\fbox{\parbox{.88\linewidth}{\centering
\textbf{One integral Seifert-form isometry}\\
Indefinite special unitary groups and strong approximation.}}\\[3pt]
$\Downarrow\quad$ simple spherical knots and weighted radial coordinates\\[3pt]
\fbox{\parbox{.88\linewidth}{\centering
\textbf{Ambiently homeomorphic hypersurface germs}\\
Disjoint sums preserve the two different ordinary orders.}}
\end{tabular}
\endgroup
\caption{The proof keeps analytic order separate from the successive
identifications of topological data. The two seed identifications are
independent; the local and global arithmetic makes them compatible.}
\label{fig:proof-route}
\end{figure}
